\scanpage{181}\mspage{36}%
% archive sheet number: 181
pouvoir se paraphraser dans le cadre des schémas excellents, à grands coups de voisinages tubulaires étales. Dans l'un et l'autre cas, le ``type d'homotopie fin'' d'un espace modéré, resp. d'un schéma excellent, se définirait par passage à la limite à partir des ``\inserted{$\unsure{\mathcal P|P}$-}types d'homotopie'' associés à des stratifications équisingulières (à strates non sing.) $\mathcal P$ de plus en plus fines. Ce type d'homotopie \inserted{``fin''} incarnerait la connaissance, non des seuls faisceaux ou \inserted{$n$-}champs loc constants, mais \inserted{(via un passage à la limite inductive)} la connaissance%
% A distinct point is at (713,1189), below the internal ai of connaissance; its origin is uncertain.
\raisebox{-1.2ex}[0pt][0pt]{\parbox{0pt}{\figure{\begin{tikzpicture}[overlay]\fill (-2.4em,0) circle[radius=0.24pt];\end{tikzpicture}}}} de \uline{tous}. Et il dépendrait, en un sens convenable, fonctoriellement%
% A distinct small point is at (700,1366), below/right of fonctoriellement, before de. Its origin and punctuation role are uncertain.
\raisebox{-0.5ex}[0pt][0pt]{\parbox{0pt}{\figure{\begin{tikzpicture}[overlay]\fill (0.2em,0) circle[radius=0.23pt];\end{tikzpicture}}}} de $X$. Dans%
% A distinct point is at (1170,1369), below the gap between Dans and le; no punctuation role is assigned.
\raisebox{-1.0ex}[0pt][0pt]{\parbox{0pt}{\figure{\begin{tikzpicture}[overlay]\fill (0.2em,0) circle[radius=0.26pt];\end{tikzpicture}}}} le cas d'un schéma de type fini sur un corps alg. clos disons, les \uline{théorèmes de finitude} cohomologique ou homotopique les plus forts s'exprimeraient précisément en termes des types d'homotopie fin, et diraient que les types d'homotopie \unsure{ordinaire} qui en sont les élément constituants sont essentiellement des ``\unsure{polyédres|polyèdres} finis'' -- et $\overline m$ des variétés à bord compactes -- ou de façon plus précise,
